%!TEX root = ../../thesis.tex
\section{背景}
\label{sec:background}

移動ロボットを目的地まで自律移動させる方法として，事前に環境の地図を作成し，走行中にその地図上で自己位置を推定しながら経路を追従する手法が広く用いられている．
この方法では，走行対象の環境ごとに地図を作成する必要があり，地図と実際の環境が食い違うと自己位置推定に失敗するおそれがある．
一方，人間は初めて訪れる建物であっても，「次の角まで直進して左に曲がる」といった道案内に従って目的地まで移動できる．
このような道案内は，正確な距離や座標を含まず，通路の形状と進む方向だけで経路を表している．

本研究グループでは，この考え方に基づき，人間が記述したシナリオとカメラ画像に基づき，移動ロボットを目的地まで移動させるシナリオ駆動型ナビゲーションを提案してきた\cite{scenario_si2023}．
この手法では，シナリオの条件と通路形状の分類結果からどの分岐路でどの方向へ進むかを決定する．
通路形状の認識と経路追従には学習ベースのモデルを用いている．
学習ベースのモデルは，学習に用いた環境では高い性能を示す一方で，外観の異なる環境では性能が低下しやすい．
そのため，走行対象の環境ごとにデータを収集して再学習することが必要になり，地図を用いないという利点が損なわれる．

近年，大規模かつ多様なデータセットで事前学習された基盤モデルが，画像認識やナビゲーションに利用されている．
いずれも学習に用いていない環境の画像に対しても有効に動作することが報告されており，環境変化への頑健性が期待できる．
各モデルについては\chapref{chap:elemental}で述べる．

ただし，ViNTや\NoMaD などの視覚ナビゲーション基盤モデルは，主に，事前に収集した画像から構築したトポロジカルグラフ上でゴール画像を指定し，そこへ移動する\cite{vint,nomad}．
そのため，基盤モデルを用いても，走行対象の環境を一度走行してゴール画像を集める必要がある．
一方，シナリオ駆動型ナビゲーションは，人間が記述した経路レベルの情報と通路形状の分類結果から目標方向を決定するため，走行対象の環境で事前にゴール画像を用意する必要がない．
そこで，両者を組み合わせ，目標方向の決定をシナリオが担い，各区間の走行を基盤モデルが担う構成とすることで，走行対象の環境での事前走行や自己位置推定を必要とせずに，目的地まで移動できることが期待される．
